\documentclass[10pt,twocolumn]{article}

\usepackage[margin=0.75in]{geometry}
\usepackage{amsmath}
\usepackage{array}
\usepackage{booktabs}
\usepackage{graphicx}
\usepackage{times}
\usepackage{url}
\usepackage[colorlinks=true,linkcolor=blue,citecolor=blue,urlcolor=blue]{hyperref}

\title{FAD-BCI: Forget-and-Adapt Deep EEG Decoding for Private Home Motor Rehabilitation}

\author{Zihao Liu\\Student ID: 251050092}
\date{}

\begin{document}

\maketitle
\begin{abstract}
Deep learning based electroencephalography (EEG) brain-computer interfaces (BCIs) are becoming plausible components of home motor rehabilitation systems, where a wearable EEG headset, an edge device, and a robotic glove or functional electrical stimulation module form a connected sensing-control loop. The central learning problem is not only to decode motor imagery accurately, but also to maintain the model as users, sessions, and consent conditions change. Cross-subject EEG models need adaptation because EEG distributions vary strongly across patients and days. Source-free domain adaptation (SFDA) is attractive because a target clinic or edge device can adapt a pretrained source model without accessing raw source EEG. However, SFDA alone does not remove source-user information already encoded in the released deep neural network. If a patient withdraws consent, a source session is later found corrupted, or a poisoned calibration session must be removed, a source-free model still needs a deletion mechanism. We propose FAD-BCI, a framework that couples machine unlearning with source-free and continual target adaptation for deep motor-imagery EEG decoding. The system first trains an EEGNet-style source decoder, then applies subject/session-level unlearning to remove selected source influence, and finally adapts the unlearned model to unlabeled target EEG streams on an edge device. The resulting model lifecycle treats forgetting, adaptation, privacy auditing, and safety gating as coupled constraints for deployable EEG rehabilitation BCIs.
\end{abstract}

\section{Introduction}
Brain-computer interfaces convert neural signals into commands, feedback, or clinical information. Motor-imagery EEG is one of the most important non-invasive BCI paradigms for rehabilitation: a patient imagines moving the left hand, right hand, feet, or another limb, and the system decodes this intention from sensorimotor rhythms. In a home rehabilitation system, the decoded intention can drive a robotic glove, functional electrical stimulation, virtual feedback, or a therapist dashboard. Such a system is naturally deployed as a connected pipeline: EEG is sensed by a wearable device, preprocessed on a phone or edge gateway, decoded by a neural model, and translated into feedback or control.

This setting makes the learning problem different from a static EEG classification benchmark. A model trained on source patients must generalize to a new patient whose neural patterns, electrode contact, scalp conductivity, fatigue level, and rehabilitation stage are different. The model must also adapt over time because EEG is non-stationary across days and sessions. At the same time, EEG is highly sensitive. It may contain subject identity, health information, cognitive state, and task-related intention. Raw EEG sharing across hospitals, cloud services, or home devices is therefore difficult to justify.

Source-free domain adaptation is a natural response to the first privacy constraint. A source institution can train a deep EEG decoder and release only the model. The target device then adapts the model to unlabeled target EEG without accessing source EEG. This is aligned with source hypothesis transfer and later black-box or source-free adaptation methods \cite{liang2020shot,liang2022dine,hwang2024sfda2,tang2025prode}. EEG-specific SFDA further shows that this learning paradigm is meaningful for cross-subject EEG and BCI settings \cite{li2025spdim,peng2025pdcc}.

However, the privacy story is incomplete. Once a deep source model is trained, the model itself may encode information about the source users. Membership inference, model inversion, and training-data reconstruction work show that trained neural models can leak training-set information through predictions, confidence scores, gradients, or internal representations \cite{fredrikson2015model,shokri2017membership,balle2022reconstructing,carlini2022membership}. In EEG, the risk is stronger because subject-specific neural patterns are not nuisance noise; they are often stable enough to support biometric or identity-related inference \cite{martinovic2012bci,frank2017subliminal,wang2025idremoval}.

This creates a lifecycle problem that SFDA alone cannot solve. Suppose a hospital trains a motor-imagery decoder using many source patients. Later, one patient withdraws consent. Or a training session is discovered to have wrong labels, electrode failure, medication-induced artifacts, or poisoning. If the model has already been released, not using source EEG in future target adaptation does not remove this source contribution from the model. The system needs a way to forget a specified user, session, feature, or label influence while preserving the remaining motor-imagery utility.

Machine unlearning addresses this deletion need. It studies how to remove the influence of selected training data from a trained model without retraining from scratch \cite{ginart2019making,guo2020certified,bourtoule2021machine,graves2021amnesiac}. Recent work investigates selective forgetting in deep networks, zero-shot unlearning, feature and label unlearning, and source-free unlearning \cite{golatkar2020eternal,chen2023zeroshot,warnecke2023features,ahmed2025sourcefree}. These ideas are directly relevant to BCI: the unit to be forgotten is often not an isolated image, but a subject, a session, or a privacy-sensitive neural feature.

We propose FAD-BCI, a \emph{forget-and-adapt} framework for deep motor-imagery EEG decoding. The framework is motivated by a concrete deployment chain. First, a source institution trains an EEGNet-style DNN on labeled source EEG. Second, before release or after a deletion request, a machine unlearning module removes selected source-user or source-session influence without relying on full retain-data retraining. Third, the unlearned model is adapted to a target user using unlabeled EEG streams on an edge device. Fourth, privacy and safety are evaluated jointly with target decoding performance.

The contributions are as follows. First, we identify a practical gap in privacy-aware BCI learning: SFDA reduces raw source-data exposure but does not address residual source-user information in the released model. Second, we formulate subject/session-level machine unlearning for deep EEG motor-imagery decoding under a source-free or retain-data-restricted constraint. Third, we design a unified objective that combines unlearning, representation privacy, source-free adaptation, uncertainty-aware pseudo-labeling, and continual test-time stability. Finally, we give an evaluation protocol that measures target utility, deletion efficacy, membership inference risk, subject-identity leakage, false command rate, and edge cost.

\section{Related Work}
\subsection{Source-Free and Black-Box Domain Adaptation}
Unsupervised domain adaptation traditionally assumes access to both source and target data. Source-free domain adaptation removes source data from the adaptation stage. SHOT adapts a source hypothesis to target data through information maximization and pseudo-labeling \cite{liang2020shot}. Model adaptation without source data, universal source-free adaptation, hypothesis transfer, DINE, and black-box domain adaptation extend this idea to more restricted settings \cite{kundu2020universal,liang2022dine}. Recent works study contrastive alignment, uncertainty, data augmentation, theoretical analysis, and proxy denoising for SFDA \cite{hwang2024sfda2,tang2025prode,mitsuzumi2024theory}. These methods motivate target-only adaptation in our framework, but they usually assume that the released source model is safe.

EEG-specific SFDA is more recent. SPDIM studies source-free unsupervised conditional and label shift adaptation in EEG \cite{li2025spdim}. Prediction consistency and confidence-based proxy domain construction are used for privacy-preserving cross-subject EEG classification \cite{peng2025pdcc}. These methods target the distribution shift problem, but they do not explicitly handle deletion of a source subject or corrupted source session from the trained model.

\subsection{Continual Test-Time Adaptation}
Home BCI deployment is not a one-shot target adaptation problem. EEG distributions drift within and across sessions. Test-time adaptation and continual test-time adaptation update a model during deployment using unlabeled target streams. Tent adapts by entropy minimization \cite{wang2021tent}. CoTTA, EATA, SAR, and related works improve stability under temporal correlation, dynamic shifts, and unreliable samples \cite{wang2022cotta,niu2022eata,niu2023sar}. We use this line of work to justify a conservative online adaptation design with confidence gating and update rejection.

\subsection{Machine Unlearning}
Machine unlearning formalizes model deletion. Early work studies data deletion, certified removal, SISA training, and amnesiac learning \cite{ginart2019making,guo2020certified,bourtoule2021machine,graves2021amnesiac}. DeltaGrad and fast unlearning methods accelerate retraining or approximate the effect of deleting training examples \cite{wu2020deltagrad}. Selective forgetting in deep networks, Bayesian class unlearning, mixed-privacy forgetting, feature-label unlearning, and instance-wise unlearning expand the space of what can be forgotten \cite{golatkar2020eternal,warnecke2023features}. Recent source-free machine unlearning is especially relevant because retain data may be unavailable or restricted \cite{ahmed2025sourcefree}. Our BCI setting uses this idea at the subject/session level.

\subsection{Privacy Leakage in Neural and EEG Models}
Membership inference and model inversion show that trained models can leak information about training data \cite{fredrikson2015model,shokri2017membership,carlini2022membership}. Training data reconstruction further demonstrates that neural networks can reveal more than aggregate task knowledge \cite{balle2022reconstructing}. In BCI, side-channel and subliminal probing attacks show that EEG responses may leak private information even when the user does not explicitly intend to disclose it \cite{martinovic2012bci,frank2017subliminal}. Recent EEG privacy works study identity removal and multisubject privacy-preserving brain decoding \cite{tian2025brainguard,wang2025idremoval}. These works motivate our use of subject-identity leakage and membership inference as evaluation metrics.

\section{Problem Setup}
\subsection{Deployment Scenario}
We consider home motor rehabilitation after stroke or motor impairment. A hospital trains a source motor-imagery EEG decoder using labeled data from multiple source patients. The deployed target patient uses a wearable EEG headset and an edge device. The edge device receives unlabeled target EEG windows and updates a small part of the model online. The output is not a free-form prediction; it is a gated rehabilitation command or feedback event.

This scenario creates three constraints. The system must personalize to a new target user without source EEG. It must forget a source user or session when consent, data quality, or security conditions change. It must avoid unsafe feedback when the model is uncertain or the EEG window is corrupted by artifacts.

\subsection{Data and Model}
Let the labeled source data be
\begin{equation}
\mathcal{D}_s=\{(x_i^s,y_i^s,u_i^s,r_i^s)\}_{i=1}^{n_s},
\end{equation}
where $x_i^s$ is a motor-imagery EEG epoch, $y_i^s$ is the movement-imagery label, $u_i^s$ is the subject ID, and $r_i^s$ is a session or recording ID. A forget request specifies a subset $\mathcal{D}_f$, such as all samples from one subject or one corrupted session. The retain set is $\mathcal{D}_r=\mathcal{D}_s\setminus \mathcal{D}_f$, but in the restricted setting $\mathcal{D}_r$ is unavailable during unlearning.

The target stream is unlabeled:
\begin{equation}
\mathcal{D}_t=\{x_j^t\}_{j=1}^{n_t}.
\end{equation}
The model is a deep EEG decoder
\begin{equation}
z=G_\theta(x), \qquad p=C_\phi(z),
\end{equation}
where $G_\theta$ is a temporal-spatial encoder and $C_\phi$ is a classifier. An EEGNet-style model is a natural default because it is compact and designed for EEG decoding \cite{lawhern2018eegnet}. The source model is trained with cross-entropy:
\begin{equation}
\mathcal{L}_{src}=\mathrm{CE}(C_\phi(G_\theta(x_s)),y_s).
\end{equation}

\subsection{Objective}
The goal is to produce a target model $M_t$ that satisfies four desiderata:
\begin{itemize}
\item \textbf{Target utility:} decode target motor imagery accurately after adaptation.
\item \textbf{Deletion efficacy:} reduce the influence of $\mathcal{D}_f$ relative to the original source model.
\item \textbf{Privacy reduction:} reduce membership inference and subject-identity leakage for forgotten data.
\item \textbf{Safety and efficiency:} support low-latency edge inference with command rejection under uncertainty.
\end{itemize}

\section{Method}
\subsection{Overview}
FAD-BCI has three stages, shown in Fig.~\ref{fig:pipeline}. Stage 1 trains a deep source EEG decoder. Stage 2 applies machine unlearning to remove a specified source subject or session. Stage 3 adapts the unlearned model to unlabeled target EEG through source-free and continual test-time objectives. The key design principle is that unlearning happens before adaptation: the target model should not inherit source influence that should have been removed.

\begin{figure}[t]
\centering
\fbox{\begin{minipage}{0.92\columnwidth}
\centering
\textbf{Source EEG training}\\
multi-subject labeled motor-imagery EEG\\
$\Downarrow$\\
EEGNet-style source decoder\\[1mm]
$\Downarrow$\\
\textbf{Machine unlearning gateway}\\
forget subject, session, feature, or label influence\\
$\Downarrow$\\
unlearned source decoder\\[1mm]
$\Downarrow$\\
\textbf{Target SFDA / CTTA}\\
unlabeled home EEG stream on edge device\\
$\Downarrow$\\
gated rehabilitation feedback
\end{minipage}}
\caption{FAD-BCI pipeline. Source-free adaptation is performed only after source-user or source-session deletion has been applied to the released decoder.}
\label{fig:pipeline}
\end{figure}

\subsection{Stage 1: Source Training}
The source institution trains $M_s=(G_{\theta_s},C_{\phi_s})$ on labeled source EEG. During this stage, it also stores model lineage metadata: subject IDs, session IDs, preprocessing configuration, random seeds, and checkpoints. These records are not used for target adaptation, but they are necessary for later deletion and audit. Without model lineage, the system cannot reliably identify what should be unlearned.

\subsection{Stage 2: Subject/Session-Level Unlearning}
The unlearning module receives $M_s$ and a forget specification $\mathcal{D}_f$. In a consent withdrawal case, $\mathcal{D}_f$ may be the last authorized deletion package for that user. In a corrupted session case, $\mathcal{D}_f$ is the invalid session. We do not assume full access to $\mathcal{D}_r$, because cross-hospital source data may be unavailable after model release.

We use an approximate unlearning objective:
\begin{equation}
\mathcal{L}_{MU}
=
\mathcal{L}_{forget}
+\lambda_a\mathcal{L}_{anchor}
+\lambda_k\mathcal{L}_{keep}
+\lambda_p\mathcal{L}_{priv}.
\end{equation}
The forgetting term reduces task memorization of the forget set:
\begin{equation}
\mathcal{L}_{forget}=-\mathrm{CE}(M_u(x_f),y_f),
\end{equation}
or uses randomized labels for $\mathcal{D}_f$. This term alone is unsafe because it can damage the model globally. Therefore, we add parameter anchoring:
\begin{equation}
\mathcal{L}_{anchor}
=\|\theta_u-\theta_s\|_2^2+\|\phi_u-\phi_s\|_2^2.
\end{equation}

The keep term preserves non-forgotten behavior without raw retain data. It can use model outputs on non-sensitive public EEG, synthetic band-limited EEG-like windows, or weakly augmented forget windows where the objective is to maintain generic feature smoothness rather than the original forget label:
\begin{equation}
\mathcal{L}_{keep}
=
\mathrm{KL}(M_s(\tilde{x}) \Vert M_u(\tilde{x})).
\end{equation}
This is not a certified guarantee; it is a practical source-free retention surrogate.

Finally, the privacy term removes subject-identifying information. A probe $A_\psi$ predicts subject or session identity from $z=G_\theta(x)$. The encoder is updated to maximize identity uncertainty on forgotten samples:
\begin{equation}
\mathcal{L}_{priv}=-H(A_\psi(G_{\theta_u}(x_f))).
\end{equation}
This connects unlearning with representation erasure and EEG identity-removal work.

\subsection{Stage 3: Source-Free Target Adaptation}
The target edge device receives the unlearned model $M_u$ and unlabeled EEG stream $\mathcal{D}_t$. The adaptation objective is
\begin{equation}
\mathcal{L}_{adapt}
=
\mathcal{L}_{MI}
+\alpha\mathcal{L}_{con}
+\beta\mathcal{L}_{unc}
+\gamma\mathcal{L}_{proto}.
\end{equation}
The mutual information term balances confident predictions and class diversity:
\begin{equation}
\mathcal{L}_{MI}
=
\frac{1}{B}\sum_{i=1}^{B} H(p_i)-H\left(\frac{1}{B}\sum_{i=1}^{B}p_i\right).
\end{equation}
Minimizing this term encourages sample-level certainty while discouraging collapse to one class. The consistency term compares predictions under EEG-preserving augmentations:
\begin{equation}
\mathcal{L}_{con}
=
\|M_t(a_1(x_t))-M_t(a_2(x_t))\|_2^2,
\end{equation}
where $a_1$ and $a_2$ may be mild temporal cropping, channel dropout, frequency masking, or Gaussian sensor noise.

The uncertainty term rejects unstable pseudo-labels. A pseudo-label is accepted only if confidence is high and predictions are stable across recent windows:
\begin{equation}
\max_k p_t(k)\geq \tau,\qquad
\mathrm{Var}(p_{t-L:t})\leq \epsilon.
\end{equation}
The prototype term maintains target class prototypes using accepted samples:
\begin{equation}
c_k \leftarrow (1-\rho)c_k+\rho z_t
\quad \text{if } \hat{y}_t=k.
\end{equation}
This makes adaptation gradual and reduces confirmation bias.

\subsection{Safety Gate}
For rehabilitation, a prediction is not automatically a command. The system emits feedback only when three checks pass: confidence, temporal stability, and artifact rejection. Otherwise, the output is ``no command'' and the sample is excluded from pseudo-label updates. This decision is important because BCI errors can cause unsafe physical feedback or misleading therapy signals.

\begin{table*}[t]
\centering
\caption{Why machine unlearning changes the design compared with SFDA-only BCI adaptation.}
\label{tab:logic}
\begin{tabular}{p{0.18\textwidth} p{0.34\textwidth} p{0.38\textwidth}}
\toprule
\textbf{Question} & \textbf{SFDA-only answer} & \textbf{FAD-BCI answer} \\
\midrule
How is raw source EEG protected? & Do not use source EEG during target adaptation. & Same, but also audit the released model as a privacy object. \\
What if a source user withdraws consent? & Not directly handled after model training. & Perform subject-level unlearning before releasing or adapting the model. \\
What if a source session is corrupted or poisoned? & Retraining may be required, or the issue is ignored. & Apply session-level unlearning and compare with oracle retraining when available. \\
Can the model still leak source membership? & Usually not evaluated. & Evaluate membership inference and subject-ID probes before and after unlearning. \\
How is target drift handled? & One-time target adaptation. & Continual source-free adaptation with uncertainty rejection. \\
\bottomrule
\end{tabular}
\end{table*}

\section{Evaluation Protocol}
\subsection{Datasets and Tasks}
The primary benchmark should be motor-imagery EEG. A minimal local setup can use PhysioNet EEGMMI runs for left/right fist and both-fists/both-feet tasks \cite{goldberger2000physionet}. A stronger evaluation can use MOABB motor-imagery datasets such as BNCI2014001 \cite{jayaram2018moabb}. The source split contains multiple subjects, the forget split is one source subject or one source session, and the target split is a held-out subject whose labels are hidden during adaptation.

\subsection{Baselines}
The evaluation should compare:
\begin{itemize}
\item \textbf{Source only:} train on source data and test on target without adaptation.
\item \textbf{SFDA only:} adapt the original source model to target EEG.
\item \textbf{MU only:} unlearn the forget set but do not adapt to the target.
\item \textbf{FAD-BCI:} unlearn first, then source-free adapt.
\item \textbf{Oracle retrain:} retrain from $\mathcal{D}_r$ and adapt. This is an upper reference, not the deployable method.
\end{itemize}

\subsection{Metrics}
Target decoding should be reported with accuracy and balanced accuracy. However, these are not enough. Unlearning efficacy should be measured by forget-set accuracy drop, distance to oracle retrain, and retain utility proxy. Privacy should be measured by membership inference AUC and subject-ID probe accuracy. Adaptation reliability should include pseudo-label acceptance rate, rejection rate, calibration error, and false command rate. Edge feasibility should include inference latency, adaptation update time, memory footprint, and raw EEG transmission volume.

\subsection{Expected Outcomes}
The method is not expected to outperform SFDA-only on pure accuracy in every case. Unlearning is a deletion and compliance mechanism, so a small accuracy cost is acceptable if privacy leakage is reduced. The desired result is that FAD-BCI approaches SFDA-only target accuracy while reducing membership inference risk and subject-ID leakage for the forgotten subject or session. If target adaptation recovers too much forgotten information, then the unlearning and adaptation objectives conflict; this conflict should be reported rather than hidden.

\section{Discussion}
\subsection{Why This Scenario Is Meaningful}
The proposed scenario is intentionally narrower than a general BCI survey. EEG motor-imagery rehabilitation is clinically meaningful, DNN decoding is justified, privacy risks are real, and source-user deletion has a clear operational meaning. This makes machine unlearning central to the problem rather than an unrelated add-on.

\subsection{Limitations}
FAD-BCI is an approximate design. Certified deletion is difficult for deep EEG models, especially when retain data are unavailable. Synthetic or public keep samples may not fully preserve retain behavior. Subject-identity probes do not capture every privacy risk. Continual adaptation can also reintroduce bias if pseudo-labels are wrong. A real deployment should combine this model-level framework with access control, encryption, consent management, secure logging, and human-in-the-loop safety.

\section{Conclusion}
This paper presented FAD-BCI, a machine unlearning assisted source-free adaptation framework for deep EEG motor-imagery decoding in home rehabilitation. The core logic is simple: target adaptation without source data is useful, but it is not sufficient when the released source model itself may retain private or invalid source-user influence. Machine unlearning supplies the missing deletion mechanism, while SFDA and continual test-time adaptation supply target personalization. Organizing the BCI model lifecycle as train, forget, adapt, and audit gives a more realistic privacy-aware design for deep BCI deployment.

\begin{thebibliography}{99}

\bibitem{lawhern2018eegnet}
Vernon J. Lawhern, Amelia J. Solon, Nicholas R. Waytowich, Stephen M. Gordon, Chou P. Hung, and Brent J. Lance. 2018. EEGNet: A Compact Convolutional Neural Network for EEG-Based Brain-Computer Interfaces. \emph{Journal of Neural Engineering} 15, 5 (2018).

\bibitem{liang2020shot}
Jian Liang, Dapeng Hu, and Jiashi Feng. 2020. Do We Really Need to Access the Source Data? Source Hypothesis Transfer for Unsupervised Domain Adaptation. In \emph{Proceedings of the 37th International Conference on Machine Learning (ICML)}.

\bibitem{kundu2020universal}
Jogendra Nath Kundu, Naveen Venkat, R. Venkatesh Babu, et al. 2020. Universal Source-Free Domain Adaptation. In \emph{Proceedings of the IEEE/CVF Conference on Computer Vision and Pattern Recognition (CVPR)}.

\bibitem{liang2022dine}
Jian Liang, Dapeng Hu, Jiashi Feng, and Ran He. 2022. DINE: Domain Adaptation from Single and Multiple Black-Box Predictors. In \emph{Proceedings of the IEEE/CVF Conference on Computer Vision and Pattern Recognition (CVPR)}.

\bibitem{hwang2024sfda2}
Uiwon Hwang, Jonghyun Lee, Juhyeon Shin, and Sungroh Yoon. 2024. SF(DA)$^2$: Source-Free Domain Adaptation Through the Lens of Data Augmentation. In \emph{International Conference on Learning Representations (ICLR)}.

\bibitem{tang2025prode}
Song Tang, Wenxin Su, Yan Gan, Mao Ye, Jianwei Zhang, and Xiatian Zhu. 2025. Proxy Denoising for Source-Free Domain Adaptation. In \emph{International Conference on Learning Representations (ICLR)}.

\bibitem{mitsuzumi2024theory}
Yu Mitsuzumi, Akisato Kimura, and Hisashi Kashima. 2024. Understanding and Improving Source-Free Domain Adaptation from a Theoretical Perspective. In \emph{Proceedings of the IEEE/CVF Conference on Computer Vision and Pattern Recognition (CVPR)}.

\bibitem{li2025spdim}
Shanglin Li, Motoaki Kawanabe, and Reinmar J. Kobler. 2025. SPDIM: Source-Free Unsupervised Conditional and Label Shift Adaptation in EEG. In \emph{International Conference on Learning Representations (ICLR)}.

\bibitem{peng2025pdcc}
Yong Peng, Jiangchuan Liu, Honggang Liu, Natasha Padfield, Junhua Li, Wanzeng Kong, Bao-Liang Lu, and Andrzej Cichocki. 2025. Prediction Consistency and Confidence-Based Proxy Domain Construction for Privacy-Preserving Cross-Subject EEG Classification. \emph{IEEE Journal of Biomedical and Health Informatics} (2025).

\bibitem{wang2021tent}
Dequan Wang, Evan Shelhamer, Shaoteng Liu, Bruno Olshausen, and Trevor Darrell. 2021. Tent: Fully Test-Time Adaptation by Entropy Minimization. In \emph{International Conference on Learning Representations (ICLR)}.

\bibitem{wang2022cotta}
Qin Wang, Olga Fink, Luc Van Gool, and Dengxin Dai. 2022. Continual Test-Time Domain Adaptation. In \emph{Proceedings of the IEEE/CVF Conference on Computer Vision and Pattern Recognition (CVPR)}.

\bibitem{niu2022eata}
Shuaicheng Niu et al. 2022. Efficient Test-Time Model Adaptation without Forgetting. In \emph{Proceedings of the 39th International Conference on Machine Learning (ICML)}.

\bibitem{niu2023sar}
Shuaicheng Niu et al. 2023. Towards Stable Test-Time Adaptation in Dynamic Wild World. In \emph{International Conference on Learning Representations (ICLR)}.

\bibitem{ginart2019making}
Antonio Ginart, Melody Y. Guan, Gregory Valiant, and James Zou. 2019. Making AI Forget You: Data Deletion in Machine Learning. In \emph{Advances in Neural Information Processing Systems (NeurIPS)}.

\bibitem{guo2020certified}
Chuan Guo, Tom Goldstein, Awni Hannun, and Laurens van der Maaten. 2020. Certified Data Removal from Machine Learning Models. In \emph{Proceedings of the 37th International Conference on Machine Learning (ICML)}.

\bibitem{bourtoule2021machine}
Lucas Bourtoule et al. 2021. Machine Unlearning. In \emph{Proceedings of the IEEE Symposium on Security and Privacy (S\&P)}.

\bibitem{graves2021amnesiac}
Laura Graves, Vineel Nagisetty, and Vijay Ganesh. 2021. Amnesiac Machine Learning. In \emph{Proceedings of the AAAI Conference on Artificial Intelligence (AAAI)}.

\bibitem{wu2020deltagrad}
Yinjun Wu, Edgar Dobriban, and Susan Davidson. 2020. DeltaGrad: Rapid Retraining of Machine Learning Models. In \emph{Proceedings of the 37th International Conference on Machine Learning (ICML)}.

\bibitem{golatkar2020eternal}
Aditya Golatkar, Alessandro Achille, and Stefano Soatto. 2020. Eternal Sunshine of the Spotless Net: Selective Forgetting in Deep Networks. In \emph{Proceedings of the European Conference on Computer Vision (ECCV)}.

\bibitem{chen2023zeroshot}
Min Chen et al. 2023. Zero-Shot Machine Unlearning. \emph{IEEE Transactions on Information Forensics and Security} (2023).

\bibitem{warnecke2023features}
Alexander Warnecke, Lukas Pirch, Christian Wressnegger, and Konrad Rieck. 2023. Machine Unlearning of Features and Labels. In \emph{Proceedings of the Network and Distributed System Security Symposium (NDSS)}.

\bibitem{ahmed2025sourcefree}
Sk Miraj Ahmed, Umit Yigit Basaran, Dripta S. Raychaudhuri, Arindam Dutta, Rohit Kundu, Fahim Faisal Niloy, Basak Guler, and Amit K. Roy-Chowdhury. 2025. Towards Source-Free Machine Unlearning. In \emph{Proceedings of the IEEE/CVF Conference on Computer Vision and Pattern Recognition (CVPR)}.

\bibitem{fredrikson2015model}
Matt Fredrikson, Somesh Jha, and Thomas Ristenpart. 2015. Model Inversion Attacks that Exploit Confidence Information and Basic Countermeasures. In \emph{Proceedings of the ACM SIGSAC Conference on Computer and Communications Security (CCS)}.

\bibitem{shokri2017membership}
Reza Shokri, Marco Stronati, Congzheng Song, and Vitaly Shmatikov. 2017. Membership Inference Attacks Against Machine Learning Models. In \emph{Proceedings of the IEEE Symposium on Security and Privacy (S\&P)}.

\bibitem{balle2022reconstructing}
Borja Balle, Giovanni Cherubin, and Jamie Hayes. 2022. Reconstructing Training Data from Trained Neural Networks. In \emph{Advances in Neural Information Processing Systems (NeurIPS)}.

\bibitem{carlini2022membership}
Nicholas Carlini et al. 2022. Membership Inference Attacks from First Principles. In \emph{Proceedings of the IEEE Symposium on Security and Privacy (S\&P)}.

\bibitem{martinovic2012bci}
Ivan Martinovic, Doug Davies, Mario Frank, Daniele Perito, Tomas Ros, and Dawn Song. 2012. On the Feasibility of Side-Channel Attacks with Brain-Computer Interfaces. In \emph{USENIX Workshop on Offensive Technologies (WOOT)}.

\bibitem{frank2017subliminal}
Mario Frank, Tiffany Hwu, and Ivan Martinovic. 2017. Subliminal Probing for Private Information via EEG-Based BCI Devices. In \emph{Proceedings on Privacy Enhancing Technologies (PETS)}.

\bibitem{tian2025brainguard}
Zhibo Tian, Ruijie Quan, Fan Ma, Kun Zhan, and Yi Yang. 2025. BrainGuard: Privacy-Preserving Multisubject Image Reconstructions from Brain Activities. In \emph{Proceedings of the AAAI Conference on Artificial Intelligence (AAAI)}.

\bibitem{wang2025idremoval}
Huabin Wang, Jie Ruan, Cunhang Fan, Yingfan Cheng, and Zhao Lv. 2025. ID-RemovalNet: Identity Removal Network for EEG Privacy Protection with Enhancing Decoding Tasks. In \emph{Proceedings of the International Joint Conference on Artificial Intelligence (IJCAI)}.

\bibitem{goldberger2000physionet}
Ary L. Goldberger et al. 2000. PhysioBank, PhysioToolkit, and PhysioNet: Components of a New Research Resource for Complex Physiologic Signals. \emph{Circulation} 101, 23 (2000), e215--e220.

\bibitem{jayaram2018moabb}
Vinay Jayaram and Alexandre Barachant. 2018. MOABB: Trustworthy Algorithm Benchmarking for BCIs. \emph{Journal of Neural Engineering} 15, 6 (2018).

\end{thebibliography}

\end{document}
